\section{Architecture A --- Clean-Sheet / Alfred-Native E2B}
\label{sec:arch-a}

Customer is beginning a new vehicle or machine architecture and is willing to adopt Alfred's preferred electrical infrastructure. Central-compute selection is out of scope; the Clicker~4 (STM32F745) stands in as the temporary central node for the entire program.

\begin{diagram}
                Clicker 4 STM32F745
                temporary central node
                         |
                    10BASE-T1S
                         |
              multidrop network
             /           |          \
            /            |           \
           v             v            v

      Alfred Node    Alfred Node   Alfred Node
       E2B mode       E2B mode      E2B mode
           |             |            |
         sensor        lighting     actuator
\end{diagram}

\subsection{What E2B precisely means}

\textbf{E2B = Ethernet-to-peripheral bridge.} It is a minimally intelligent edge translation node: physical peripheral interface $\leftrightarrow$ standardized Ethernet representation, and nothing upstream of that. Responsibilities: receive T1S data; translate commands into PWM/GPIO/UART/SPI; acquire peripheral state; sample ADC inputs; capture sensor data; publish state over T1S; report basic node health; handle the minimum local deterministic hardware function a peripheral genuinely requires (e.g., a PWM period, not a control law). An E2B is \emph{not} a zonal computer, an application ECU, or a vehicle-behavior controller -- that logic stays upstream, on whatever plays the central-compute role (Clicker~4 for this program).

\begin{diagram}
                  10BASE-T1S
                       |
                       v
               +---------------+
               | Alfred Node   |
               | E2B config    |
               +-------+-------+
                       |
         +-------------+-------------+
         |             |             |
        PWM           ADC           SPI
         |             |             |
      actuator       sensor      peripheral
\end{diagram}

Different E2B applications populate different peripheral-interface footprints on the common node (Section~\ref{sec:common-platform}) rather than requiring a different board design; a daughtercard/interface-module approach is not needed at this node count and is explicitly deferred (Section~\ref{sec:feature-cuts}) unless Rev~C schedule margin allows a connectorized peripheral header as a stretch goal.

\subsection{10BASE-T1S architecture assumptions}

\begin{center}
\small
\begin{tabularx}{\linewidth}{@{}p{3.0cm} X@{}}
\toprule
\textbf{Property} & \textbf{Prototype assumption for this program} \\
\midrule
Topology & Multidrop, unshielded twisted pair, bus (not star); short bench/rig runs (a few meters), not the full 25\,m/8-node IEEE~802.3cg reference topology -- that headroom is not needed to prove the architecture in September--November. \\
Medium access & PLCA (Physical Layer Collision Avoidance); CSMA/CD fallback exists in the standard but is not the validation target -- PLCA is. \\
Node addressing & Node~0 = PLCA coordinator (the Clicker~4 side, or a dedicated coordinator node); followers numbered 1--N; \emph{transmit opportunities} register set to (max node ID + 1) per LAN8650/1 configuration app note. \\
Node count & 2 nodes minimum for the November pass gate, 3 preferred (Section~\ref{sec:arch-a-demo}); LAN8650/1 supports up to 9 transmit opportunities (practically $\geq$8 followers), far beyond what this program needs to validate. \\
Cable/termination & Standard twisted pair (e.g., Cat5e-class stranded pair is sufficient for bench use); $100\,\Omega$ end-of-segment termination via the MikroE click-board jumpers (both populated at each physical end, unpopulated at interior drop points). \\
Coupling & Transformer-coupled per Microchip reference design (no direct-DC bench shortcuts); connectors: screw-terminal or keyed 2-pin for bench/Rev~A/B, moving to a locking automotive-style connector only at Rev~C. \\
Startup / joining & Coordinator beacons; followers synchronize to the beacon and take their assigned transmit opportunity; joining is passive from the follower's perspective (no discovery handshake beyond PLCA's own beacon mechanism at this layer). \\
Loss / recovery & A dropped follower is simply silent for its opportunity slot (no bus lockup); recovery on reconnect is measured, not assumed (Section~\ref{sec:validation-plan}). \\
Bandwidth & 10\,Mbps shared, half-duplex; realistic sustained per-node throughput is a small fraction of that once framing/opportunity overhead and node count are accounted for -- sufficient for command/telemetry traffic, explicitly not sized for bulk sensor streaming. \\
\bottomrule
\end{tabularx}
\end{center}

\begin{decision}[title={Why 10BASE-T1S is the wedge for a clean-sheet customer}]
It offers Ethernet-native communication end to end (same frame format the customer's tooling, logging, and eventual central compute already understand), multidrop topology on a single twisted pair (fewer point-to-point runs than a star), fewer protocol-specific ECU implementations (one physical/link layer instead of a CAN stack plus a LIN stack plus discrete GPIO harnessing), and straightforward expansion (add a node, do not redesign a bus segment's arbitration budget). It does not solve every network requirement -- it is not a substitute for 100BASE-T1/1000BASE-T1 where bandwidth is genuinely needed, and it does not by itself provide the semantic/application layer; it lowers the cost of getting a peripheral onto the network in the first place, which is the actual sales argument.
\end{decision}

\subsection{Greenfield November demonstration}
\label{sec:arch-a-demo}

\begin{diagram}
                 Clicker 4
                     |
                10BASE-T1S
                     |
          +----------+----------+----------+
          |                     |          |
          v                     v          v

      Node Rev C            Node Rev C   Node Rev C
       E2B mode              E2B mode     E2B mode
          |                     |            |
        Sensor               Actuator     (stretch: 3rd role)
\end{diagram}

Minimum accepted for the November~30 gate: Clicker~4 plus two Alfred nodes in E2B mode (one sensor-facing, one actuator-facing) on a multidrop 10BASE-T1S segment; three nodes preferred if schedule allows. Demonstrated live: network bring-up, multidrop operation with $\geq$2 simultaneous followers, sensor-to-Clicker data flow, Clicker-to-actuator command flow, node disconnect, node reconnect, node power-cycle, and PLCA-level error recovery. Measured and recorded per the procedures in Section~\ref{sec:validation-plan}: end-to-end and node-to-node latency, jitter, bus utilization, error rate, cold-start time, reconnect time, and per-node supply current (idle and active).
